
\title{\textbf{Mathematical Foundations, Architecture Diagrams, and State Analysis\\
of Modern LLM Training Methods --- Part 2:\\
Preference Alignment, Reinforcement Learning, and Reward Modeling}}
\author{Chief Strategist J}
\date{\today}
\maketitle
\tableofcontents
\newpage

%% ============================================================
\part{Preference Alignment Methods}
%% ============================================================

\section{Reinforcement Learning from Human Feedback (RLHF)}

\subsection{Mathematical Foundation}
RLHF aligns language models with human preferences in a two-stage process. 

First, a Reward Model (RM) $r_\phi(x, y)$ is trained on a dataset of human comparisons $\mathcal{D} = \{(x, y_w, y_l)\}$, where $y_w$ is preferred to $y_l$ given prompt $x$. The preference is modeled using the Bradley-Terry preference model:
\begin{equation}
  P(y_w \succ y_l \mid x) = \sigma\bigl( r_\phi(x, y_w) - r_\phi(x, y_l) \bigr) = \frac{1}{1 + \exp\bigl( - (r_\phi(x, y_w) - r_\phi(x, y_l)) \bigr)}
\end{equation}
The RM is trained by minimizing the negative log-likelihood of this pairwise choice:
\begin{equation}
  \mathcal{L}_{\text{RM}}(\phi) = -\mathbb{E}_{(x, y_w, y_l) \sim \mathcal{D}} \left[ \log \sigma \bigl( r_\phi(x, y_w) - r_\phi(x, y_l) \bigr) \right]
\end{equation}

Second, the SFT policy parameters $\theta$ are fine-tuned to maximize the expected reward under the learned RM, while regularized by a Kullback-Leibler (KL) divergence constraint to the reference policy $\pi_{\text{ref}}$ (preventing policy collapse and reward hacking):
\begin{equation}
  \mathcal{J}_{\text{RLHF}}(\theta) = \mathbb{E}_{x \sim \mathcal{D}_x, y \sim \pi_\theta(\cdot \mid x)} \left[ r_\phi(x, y) - \beta D_{\text{KL}}\bigl( \pi_\theta(y \mid x) \parallel \pi_{\text{ref}}(y \mid x) \bigr) \right]
\end{equation}
where $\beta$ is a regularizing coefficient, and the KL term at each token generation step is factorized as:
\begin{equation}
  D_{\text{KL}}\bigl( \pi_\theta(y \mid x) \parallel \pi_{\text{ref}}(y \mid x) \bigr) = \sum_{t=1}^T \log \frac{\pi_\theta(y_t \mid x, y_{<t})}{\pi_{\text{ref}}(y_t \mid x, y_{<t})}
\end{equation}

\subsection{Architecture Flow Diagram}

\begin{figure}[h]
\centering
\begin{adjustbox}{max width=\linewidth}
\begin{tikzpicture}[node distance=0.7cm and 1.1cm]
  % Phase 1: RM training
  \node[datblk] (human) {\textbf{Human} \\ \textbf{Comparisons} \\ ($(y_w \succ y_l)$)};
  \node[trnblk, right=of human] (rm) {\textbf{Reward} \\ (Model $r_\phi$)};

  % Phase 2: PPO
  \node[frzblk, below=1.8cm of human] (ref) {\textbf{Reference} \\ (Policy $\pi_{\text{ref}}$)};
  \node[trnblk, right=of ref] (policy) {\textbf{Policy} \\ ($\pi_\theta$ (PPO))};
  \node[grayblk, right=of policy] (sample) {\textbf{Sample} \\ ($y \sim \pi_\theta$)};
  \node[lossblk, right=of sample] (rloss) {\textbf{$r_\phi(x,y)$} \\ ($-\beta$KL)};
  \node[mrgblk, below=0.8cm of rloss] (aligned) {\textbf{Aligned} \\ (Model $\pi^*$)};

  \draw[arr] (human)--(rm);
  \draw[arr] (rm.south) -- ++(0,-0.4) -| (rloss.north);
  \draw[arr] (ref)--(policy);
  \draw[arr] (policy)--(sample);
  \draw[arr] (sample)--(rloss);
  \draw[arr] (rloss)--(aligned);
  \draw[arr,dashed,blue!60] (policy.north) -- ++(0,0.4) -| (ref.north) node[midway,above,font=\tiny]{KL penalty};

  \node[font=\small\bfseries,red] at ($(human)!0.5!(rm)+(0,0.8)$) {Phase 1: RM};
  \node[font=\small\bfseries,blue!70] at ($(ref)!0.5!(policy)+(0,-0.7)$) {Phase 2: PPO};
\end{tikzpicture}
\end{adjustbox}
\caption{RLHF: reward model trained on human comparisons then used to guide PPO policy optimisation.}
\end{figure}

\subsection{State Transition Table}
\begin{table}[h]\centering
\begin{tabular}{lll}\toprule
\textbf{Property} & \textbf{Before} & \textbf{After}\\\midrule
Policy type & SFT checkpoint & RLHF-aligned policy\\
Reward signal & None & Human preference score\\
KL from reference & 0 & Bounded by $\beta$\\
Human preference win-rate & $\sim$50\% & $>70$\%\\
Training complexity & Single loop & Two-phase (RM + PPO)\\\bottomrule
\end{tabular}
\caption{State properties before and after RLHF.}
\end{table}

\newpage
%% ------------------------------------------------------------
\section{Direct Preference Optimization (DPO)}

\subsection{Mathematical Foundation}
DPO eliminates the need for training an explicit reward model or running reinforcement learning. Starting from the KL-constrained RL objective:
\begin{equation}
  \max_{\pi} \mathbb{E}_{x, y \sim \pi} \left[ r(x, y) \right] - \beta D_{\text{KL}}(\pi(y \mid x) \parallel \pi_{\text{ref}}(y \mid x))
\end{equation}
we formulate the unconstrained objective using a partition function $Z(x)$:
\begin{equation}
  \max_{\pi} \mathbb{E}_{x} \left[ \mathbb{E}_{y \sim \pi} \left[ r(x, y) - \beta \log \frac{\pi(y \mid x)}{\pi_{\text{ref}}(y \mid x)} \right] \right]
\end{equation}
The closed-form optimal policy solution $\pi^*(y \mid x)$ under this objective is:
\begin{equation}
  \pi^*(y \mid x) = \frac{1}{Z(x)} \pi_{\text{ref}}(y \mid x) \exp\left( \frac{r(x, y)}{\beta} \right), \quad Z(x) = \sum_{y} \pi_{\text{ref}}(y \mid x) \exp\left( \frac{r(x, y)}{\beta} \right)
\end{equation}
By rearranging this relation, we express the implicit reward function $r(x, y)$ in terms of the policy likelihoods:
\begin{equation}
  r(x, y) = \beta \log \frac{\pi^*(y \mid x)}{\pi_{\text{ref}}(y \mid x)} + \beta \log Z(x)
\end{equation}
Substituting this expression into the Bradley-Terry preference probability $P(y_w \succ y_l \mid x) = \sigma(r(x, y_w) - r(x, y_l))$, the partition function $Z(x)$ cancels out, resulting in the DPO loss function:
\begin{equation}
  \mathcal{L}_{\text{DPO}}(\theta; \pi_{\text{ref}}) = -\mathbb{E}_{(x, y_w, y_l) \sim \mathcal{D}} \left[ \log \sigma \left( \beta \log \frac{\pi_\theta(y_w \mid x)}{\pi_{\text{ref}}(y_w \mid x)} - \beta \log \frac{\pi_\theta(y_l \mid x)}{\pi_{\text{ref}}(y_l \mid x)} \right) \right]
\end{equation}

\subsection{Architecture Flow Diagram}

\begin{figure}[h]
\centering
\begin{adjustbox}{max width=\linewidth}
\begin{tikzpicture}[node distance=0.7cm and 1.1cm]
  \node[datblk] (pairs) {\textbf{Preference Pairs} \\ ($(x, y_w, y_l)$)};
  \node[trnblk, above right=0.6cm and 1.5cm of pairs] (pol) {\textbf{Policy} \\ ($\pi_\theta$)};
  \node[frzblk, below right=0.6cm and 1.5cm of pairs] (ref) {\textbf{Reference} \\ ($\pi_{\text{ref}}$)};
  \node[lossblk, right=3.0cm of pairs] (ratio) {\textbf{Log Ratio} \\ ($\log\frac{\pi_\theta}{\pi_{\text{ref}}}$)};
  \node[mrgblk, right=of ratio] (out) {\textbf{DPO} \\ (Loss)};

  \draw[arr] (pairs.east) |- (pol.west);
  \draw[arr] (pairs.east) |- (ref.west);
  \draw[arr] (pol.east) -| (ratio.north);
  \draw[arr] (ref.east) -| (ratio.south);
  \draw[arr] (ratio)--(out);

  \node[font=\tiny,red,align=center,below=0.3cm of ratio] {no RM needed\\no sampling loop};
\end{tikzpicture}
\end{adjustbox}
\caption{DPO: both policy and frozen reference score chosen/rejected responses; loss uses their log-ratio difference.}
\end{figure}

\subsection{State Transition Table}
\begin{table}[h]\centering
\begin{tabular}{lll}\toprule
\textbf{Property} & \textbf{Before (RLHF)} & \textbf{After (DPO)}\\\midrule
Reward model & Explicit $r_\phi$ & Implicit in log-ratio\\
Sampling loop & Yes (PPO rollouts) & No\\
Reference policy & Frozen copy & Frozen copy\\
Training stability & Sensitive to reward hacking & More stable\\
Compute cost & High (2-phase) & Lower (single SFT-like pass)\\\bottomrule
\end{tabular}
\caption{State properties comparing RLHF and DPO.}
\end{table}

\newpage
%% ------------------------------------------------------------
\section{Identity Preference Optimization (IPO)}

\subsection{Mathematical Foundation}
IPO relaxes the Bradley-Terry preference assumption of DPO to prevent overfitting and improve generalization on preference pairs. By applying identity mapping directly to the pairwise probability, the IPO objective minimizes the mean squared error between the policy log-ratio difference and a target preference margin:
\begin{equation}
  \mathcal{L}_{\text{IPO}}(\theta) = -\mathbb{E}_{(x, y_w, y_l) \sim \mathcal{D}} \left[ \log \sigma \left( \beta \log \frac{\pi_\theta(y_w \mid x)}{\pi_{\text{ref}}(y_w \mid x)} - \beta \log \frac{\pi_\theta(y_l \mid x)}{\pi_{\text{ref}}(y_l \mid x)} \right) \right] + \text{regularization}
\end{equation}
Formally, IPO minimizes the quadratic loss of the difference of log-ratios relative to a regularizing parameter $\tau$:
\begin{equation}
  \mathcal{L}_{\text{IPO}}(\theta) = \mathbb{E}_{(x, y_w, y_l) \sim \mathcal{D}} \left[ \left( \log \frac{\pi_\theta(y_w \mid x)}{\pi_{\text{ref}}(y_w \mid x)} - \log \frac{\pi_\theta(y_l \mid x)}{\pi_{\text{ref}}(y_l \mid x)} - \frac{\tau}{2 \beta} \right)^2 \right]
\end{equation}
This bounds the likelihood shift of both preferred and dispreferred sequences, preventing the model from outputting degenerate probabilities for preferred responses.

\subsection{Architecture Flow Diagram}

\begin{figure}[h]
\centering
\begin{adjustbox}{max width=\linewidth}
\begin{tikzpicture}[node distance=0.7cm and 1.1cm]
  \node[datblk] (pairs) {\textbf{Pairs} \\ ($(x,y_w,y_l)$)};
  \node[trnblk, above right=0.5cm and 1.4cm of pairs] (pol) {\textbf{Policy $\pi_\theta$}};
  \node[frzblk, below right=0.5cm and 1.4cm of pairs] (ref) {\textbf{Ref $\pi_{\text{ref}}$}};
  \node[grayblk, right=3.2cm of pairs] (h) {\textbf{Gap $h$} \\ (log-ratio diff)};
  \node[lossblk, right=of h] (loss) {\textbf{Squared} \\ (Loss $\mathcal{L}_{\text{IPO}}$)};
  \node[mrgblk, right=of loss] (tgt) {\textbf{Target} \\ ($\frac{1}{2\beta}$)};

  \draw[arr] (pairs.east)|-  (pol.west);
  \draw[arr] (pairs.east)|-  (ref.west);
  \draw[arr] (pol.east) -| (h.north);
  \draw[arr] (ref.east) -| (h.south);
  \draw[arr] (h)--(loss);
  \draw[arr] (tgt.west)--(loss.east);

  \node[font=\tiny,align=center,below=0.3cm of loss] {bounded: $h \to \frac{1}{2\beta}$};
\end{tikzpicture}
\end{adjustbox}
\caption{IPO: like DPO but the log-ratio gap is penalised toward a finite target $1/(2\beta)$.}
\end{figure}

\subsection{State Transition Table}
\begin{table}[h]\centering
\begin{tabular}{lll}\toprule
\textbf{Property} & \textbf{Before (DPO)} & \textbf{After (IPO)}\\\midrule
Loss type & Log-sigmoid & Squared error\\
Divergence at optimum & Unbounded & Bounded by $\frac{1}{2\beta}$\\
Target preference gap & Maximise & Converge to $\frac{1}{2\beta}$\\
Regularisation & KL implicit & Explicit squared bound\\
Training saturation & Can saturate & More stable\\\bottomrule
\end{tabular}
\caption{State properties comparing DPO and IPO.}
\end{table}

\newpage
%% ------------------------------------------------------------
\section{Odds Ratio Preference Optimization (ORPO)}

\subsection{Mathematical Foundation}
ORPO combines SFT loss and an odds-ratio contrast in a single objective — no reference model needed. The odds of generating $y$ under policy $\pi_\theta$ is $\text{odds}(y|x) = \pi_\theta(y|x)/(1-\pi_\theta(y|x))$. The combined loss is:
\begin{equation}
  \mathcal{L}_{\text{ORPO}} = \mathcal{L}_{\text{SFT}} - \lambda\,\mathbb{E}\!\left[\log\sigma\!\left(\log\frac{\text{odds}(y_w|x)}{\text{odds}(y_l|x)}\right)\right]
\end{equation}
The SFT term increases $\pi_\theta(y_w|x)$ while the odds-ratio term simultaneously penalises $y_l$ without requiring a frozen reference copy.

\subsection{Architecture Flow Diagram}

\begin{figure}[h]
\centering
\begin{adjustbox}{max width=\linewidth}
\begin{tikzpicture}[node distance=0.7cm and 1.0cm]
  \node[datblk] (data) {\textbf{$(x, y_w, y_l)$}};
  \node[trnblk, right=of data] (pol) {\textbf{Policy} \\ ($\pi_\theta$)};
  \node[lossblk, above right=0.5cm and 1.4cm of pol] (sft) {\textbf{SFT Loss} \\ (on $y_w$)};
  \node[lossblk, below right=0.5cm and 1.4cm of pol] (or) {\textbf{Odds Ratio} \\ (Contrast)};
  \node[mrgblk, right=3.2cm of pol] (total) {\textbf{ORPO} \\ (Total Loss)};

  \draw[arr] (data)--(pol);
  \draw[arr] (pol.east) |- (sft.west);
  \draw[arr] (pol.east) |- (or.west);
  \draw[arr] (sft.east) -| (total.north);
  \draw[arr] (or.east) -| (total.south);

  \node[font=\tiny,red,align=center,below=0.2cm of or] {no ref model};
\end{tikzpicture}
\end{adjustbox}
\caption{ORPO: SFT loss on the chosen response combined with an odds-ratio preference contrast.}
\end{figure}

\subsection{State Transition Table}
\begin{table}[h]\centering
\begin{tabular}{lll}\toprule
\textbf{Property} & \textbf{Before} & \textbf{After}\\\midrule
Reference model & Required (DPO/RLHF) & Not required\\
Training phases & 2 (SFT then align) & 1 (joint)\\
Loss terms & Single & SFT + odds-ratio\\
Memory overhead & $2\times$ model & $1\times$ model only\\
Data format & Prompt + completion & Triplet $(x, y_w, y_l)$\\\bottomrule
\end{tabular}
\caption{State properties before and after ORPO.}
\end{table}

\newpage
%% ------------------------------------------------------------
\section{Kahneman-Tversky Optimization (KTO)}

\subsection{Mathematical Foundation}
KTO aligns language models using binary preference labels (e.g., correct/incorrect, helpful/harmful) instead of comparative pairs. This objective mathematically models human utility based on Kahneman and Tversky's prospect theory, which describes how humans weigh losses and gains asymmetricaly.
Let $y$ be a response to prompt $x$, and $r \in \{0, 1\}$ be its correctness label. The utility of the policy is defined as:
\begin{equation}
  \mathcal{L}_{\text{KTO}}(\pi_\theta; \pi_{\text{ref}}) = -\mathbb{E}_{x, y} \left[ w(y) \cdot v\left( \log \frac{\pi_\theta(y \mid x)}{\pi_{\text{ref}}(y \mid x)} - z(x) \right) \right]
\end{equation}
where $v(u)$ is the value function mapping gains and losses:
\begin{equation}
  v(u) = \begin{cases} u & \text{if } u > 0 \\ \lambda u & \text{if } u \le 0 \end{cases}
\end{equation}
where $\lambda > 1$ represents loss aversion. The term $z(x)$ is a running reference point representing the expected log-ratio of the policy relative to the reference:
\begin{equation}
  z(x) = \mathbb{E}_{y' \sim \pi_{\text{ref}}(\cdot \mid x)} \left[ \log \frac{\pi_\theta(y' \mid x)}{\pi_{\text{ref}}(y' \mid x)} \right]
\end{equation}
and $w(y)$ is a weighting factor balancing positive and negative feedback frequencies.

\subsection{Architecture Flow Diagram}

\begin{figure}[h]
\centering
\begin{adjustbox}{max width=\linewidth}
\begin{tikzpicture}[node distance=0.7cm and 1.0cm]
  \node[datblk] (data) {\textbf{Individual Samples} \\ ($(x, y, z\in\{0,1\})$)};
  \node[trnblk, right=of data] (pol) {\textbf{Policy $\pi_\theta$}};
  \node[frzblk, below=0.8cm of pol] (ref) {\textbf{Ref $\pi_{\text{ref}}$}};
  \node[grayblk, right=of pol] (r) {\textbf{Implicit} \\ (Reward $r$)};
  \node[grayblk, right=of ref] (kl) {\textbf{KL} \\ (Baseline)};
  \node[lossblk, right=2.0cm of r] (util) {\textbf{Utility} \\ ($v(y,z)$)};
  \node[mrgblk, right=of util] (out) {\textbf{KTO} \\ (Loss)};

  \draw[arr] (data)--(pol);
  \draw[arr] (data.south) -- ++(0,-0.3) -| (ref.west);
  \draw[arr] (pol)--(r);
  \draw[arr] (ref)--(kl);
  \draw[arr] (r.east) -| (util.north);
  \draw[arr] (kl.east) -| (util.south);
  \draw[arr] (util)--(out);
\end{tikzpicture}
\end{adjustbox}
\caption{KTO: each sample is individually labelled good/bad; utility function uses prospect theory asymmetry.}
\end{figure}

\subsection{State Transition Table}
\begin{table}[h]\centering
\begin{tabular}{lll}\toprule
\textbf{Property} & \textbf{Before} & \textbf{After}\\\midrule
Data format & Pairwise $(y_w, y_l)$ & Individual $(y, z)$\\
Data collection effort & High (pairwise) & Lower (binary label)\\
Loss-aversion modelling & None & Asymmetric $w_+, w_-$\\
Reference model & Optional & Used for KL baseline\\
Applicable at scale & Full pairs required & Works with any ratio\\\bottomrule
\end{tabular}
\caption{State properties before and after KTO.}
\end{table}

\newpage
%% ------------------------------------------------------------
\section{Simple Preference Optimization (SimPO)}

\subsection{Mathematical Foundation}
SimPO improves upon DPO by removing the reference policy $\pi_{\text{ref}}$ during training, reducing memory usage, and introducing a length-normalized reward with a target margin $\gamma$.
The SimPO implicit reward for a sequence $y$ given $x$ is defined as the average log-likelihood:
\begin{equation}
  r_{\text{SimPO}}(x, y) = \frac{\beta}{|y|} \log \pi_\theta(y \mid x)
\end{equation}
where $|y|$ is the sequence length in tokens. The SimPO pairwise loss function is formulated as:
\begin{equation}
  \mathcal{L}_{\text{SimPO}}(\theta) = -\mathbb{E}_{(x, y_w, y_l) \sim \mathcal{D}} \left[ \log \sigma \left( \frac{\beta}{|y_w|} \log \pi_\theta(y_w \mid x) - \frac{\beta}{|y_l|} \log \pi_\theta(y_l \mid x) - \gamma \right) \right]
\end{equation}
where $\gamma > 0$ is a target margin that forces the average likelihood of the winning response to exceed the losing response by a strict buffer, preventing over-optimization of short sequences.

\subsection{Architecture Flow Diagram}

\begin{figure}[h]
\centering
\begin{adjustbox}{max width=\linewidth}
\begin{tikzpicture}[node distance=0.7cm and 1.0cm]
  \node[datblk] (data) {\textbf{$(x, y_w, y_l)$}};
  \node[trnblk, right=of data] (pol) {\textbf{Policy $\pi_\theta$} \\ (only copy)};
  \node[grayblk, above right=0.5cm and 1.2cm of pol] (rw) {\textbf{$r_w = \frac{\beta}{|y_w|}\log\pi(y_w|x)$}};
  \node[grayblk, below right=0.5cm and 1.2cm of pol] (rl) {\textbf{$r_l = \frac{\beta}{|y_l|}\log\pi(y_l|x)$}};
  \node[lossblk, right=3.6cm of pol] (loss) {\textbf{SimPO} \\ (Loss $-\log\sigma(r_w - r_l - \gamma)$)};

  \draw[arr] (data)--(pol);
  \draw[arr] (pol.east) |- (rw.west);
  \draw[arr] (pol.east) |- (rl.west);
  \draw[arr] (rw.east) -| (loss.north);
  \draw[arr] (rl.east) -| (loss.south);

  \node[red,font=\tiny,align=center,below=0.3cm of pol] {no reference\\model needed};
\end{tikzpicture}
\end{adjustbox}
\caption{SimPO: length-normalised log-probs replace the reference log-ratio; margin $\gamma$ enforces quality gap.}
\end{figure}

\subsection{State Transition Table}
\begin{table}[h]\centering
\begin{tabular}{lll}\toprule
\textbf{Property} & \textbf{Before (DPO)} & \textbf{After (SimPO)}\\\midrule
Reference model & Required & Not needed\\
Reward definition & Log-ratio $\pi/\pi_{\text{ref}}$ & Length-norm avg log-prob\\
Length bias & Possible & Mitigated by $1/|y|$\\
Margin hyperparameter & None & $\gamma > 0$\\
Memory saving & None & No second model forward pass\\\bottomrule
\end{tabular}
\caption{State properties comparing DPO and SimPO.}
\end{table}

\newpage
%% ------------------------------------------------------------
\section{Constrained Preference Optimization (CPO)}

\subsection{Mathematical Foundation}
CPO formulates alignment as a constrained optimisation: maximise reward while staying within a trust region $\epsilon$ of the reference policy measured in reverse KL divergence:
\begin{equation}
  \max_{\pi_\theta} \mathbb{E}[r(x,y)] \quad \text{s.t.} \quad \mathbb{KL}[\pi_{\text{ref}} \| \pi_\theta] \leq \epsilon
\end{equation}
The Lagrangian relaxation yields a combined training objective:
\begin{equation}
  \mathcal{L}_{\text{CPO}} = -\mathbb{E}[r(x,y)] + \mu\bigl(\mathbb{KL}[\pi_{\text{ref}} \| \pi_\theta] - \epsilon\bigr)
\end{equation}
where $\mu \geq 0$ is a dual variable updated online to enforce the constraint.

\subsection{Architecture Flow Diagram}

\begin{figure}[h]
\centering
\begin{adjustbox}{max width=\linewidth}
\begin{tikzpicture}[node distance=0.7cm and 1.1cm]
  \node[datblk] (data) {\textbf{$(x, y_w, y_l)$}};
  \node[trnblk, right=of data] (pol) {\textbf{Policy $\pi_\theta$}};
  \node[frzblk, below=0.9cm of pol] (ref) {\textbf{Ref $\pi_{\text{ref}}$}};
  \node[lossblk, right=of pol] (rew) {\textbf{Reward} \\ (Objective)};
  \node[grayblk, right=of ref] (kl) {\textbf{KL Constraint} \\ ($\leq \epsilon$)};
  \node[orangeblk, right=2.5cm of pol] (mu) {\textbf{Dual Var $\mu$} \\ (update online)};
  \node[mrgblk, below=0.7cm of mu] (out) {\textbf{CPO Loss}};

  \draw[arr] (data)--(pol);
  \draw[arr] (pol)--(rew);
  \draw[arr] (pol.south)--(ref.north);
  \draw[arr] (ref)--(kl);
  \draw[arr] (rew.east)-|(mu.north);
  \draw[arr] (kl.east)-|(mu.south);
  \draw[arr] (mu)--(out);
\end{tikzpicture}
\end{adjustbox}
\caption{CPO: reward maximisation with an explicit KL constraint; dual variable $\mu$ adapts online.}
\end{figure}

\subsection{State Transition Table}
\begin{table}[h]\centering
\begin{tabular}{lll}\toprule
\textbf{Property} & \textbf{Before} & \textbf{After}\\\midrule
Constraint form & Soft KL penalty & Hard KL constraint $\leq \epsilon$\\
Dual variable $\mu$ & Fixed & Dynamically updated\\
Policy deviation & Unbounded & Bounded within trust region\\
Reward maximisation & Soft & Subject to constraint\\
Training dynamics & Standard & Saddle-point optimisation\\\bottomrule
\end{tabular}
\caption{State properties before and after CPO.}
\end{table}

\newpage
%% ------------------------------------------------------------
\section{Anchored Preference Optimization (APO)}

\subsection{Mathematical Foundation}
APO introduces an anchor term that explicitly regularises the policy toward the reference model while simultaneously pushing chosen responses up and rejected responses down:
\begin{equation}
  \mathcal{L}_{\text{APO}} = \mathbb{E}\!\left[\underbrace{-\log\sigma\bigl(\beta(r_w - r_l)\bigr)}_{\text{preference}} + \lambda\underbrace{\bigl(\log\frac{\pi_\theta(y_w|x)}{\pi_{\text{ref}}(y_w|x)} - \log\frac{\pi_\theta(y_l|x)}{\pi_{\text{ref}}(y_l|x)}\bigr)^2}_{\text{anchor}}\right]
\end{equation}
The anchor term is minimised when both the chosen and rejected log-ratios equal a fixed reference, preventing the model from drifting.

\subsection{Architecture Flow Diagram}

\begin{figure}[h]
\centering
\begin{adjustbox}{max width=\linewidth}
\begin{tikzpicture}[node distance=0.7cm and 1.0cm]
  \node[datblk] (data) {\textbf{$(x,y_w,y_l)$}};
  \node[trnblk, right=of data] (pol) {\textbf{$\pi_\theta$}};
  \node[frzblk, below=0.9cm of pol] (ref) {\textbf{$\pi_{\text{ref}}$}};
  \node[lossblk, above right=0.5cm and 1.4cm of pol] (pref) {\textbf{Preference} \\ (Term)};
  \node[orangeblk, below right=0.5cm and 1.4cm of pol] (anchor) {\textbf{Anchor} \\ (Regulariser)};
  \node[mrgblk, right=3.3cm of pol] (total) {\textbf{APO} \\ (Loss)};

  \draw[arr] (data)--(pol);
  \draw[arr] (pol.east) |- (pref.west);
  \draw[arr] (pol.south)--(ref.north);
  \draw[arr] (ref.east) -| (anchor.south);
  \draw[arr] (pol.east) |- (anchor.west);
  \draw[arr] (pref.east) -| (total.north);
  \draw[arr] (anchor.east) -| (total.south);
\end{tikzpicture}
\end{adjustbox}
\caption{APO: combines standard preference loss with an anchor regulariser against the reference model.}
\end{figure}

\subsection{State Transition Table}
\begin{table}[h]\centering
\begin{tabular}{lll}\toprule
\textbf{Property} & \textbf{Before (DPO)} & \textbf{After (APO)}\\\midrule
Anchor regularisation & Implicit KL & Explicit squared anchor\\
Reference role & Log-ratio denominator & Explicit target\\
Policy drift & Can drift & Penalised\\
Loss terms & 1 & 2 (preference + anchor)\\
Stability & Moderate & Higher (dual grounding)\\\bottomrule
\end{tabular}
\caption{State properties comparing DPO and APO.}
\end{table}

\newpage
%% ------------------------------------------------------------
\section{Rank Responses from Human Feedback (RRHF)}

\subsection{Mathematical Foundation}
RRHF generates $K$ candidate responses $\{y^{(1)}, \ldots, y^{(K)}\}$ from the model, ranks them by human or AI annotators, and trains the model to assign higher log-probability to higher-ranked responses via a pairwise ranking loss:
\begin{equation}
  \mathcal{L}_{\text{RRHF}} = \sum_{i < j}\max\!\bigl(0,\, \text{score}(y^{(j)}) - \text{score}(y^{(i)}) + \delta\bigr) + \lambda \mathcal{L}_{\text{SFT}}(y^{(1)})
\end{equation}
where $\text{score}(y^{(k)}) = \log P_\theta(y^{(k)}|x)/|y^{(k)}|$ is the normalised log-prob. The SFT term anchors the top-ranked response.

\subsection{Architecture Flow Diagram}

\begin{figure}[h]
\centering
\begin{adjustbox}{max width=\linewidth}
\begin{tikzpicture}[node distance=0.6cm and 1.0cm]
  \node[datblk] (x) {\textbf{Prompt $x$}};
  \node[trnblk, right=of x] (gen) {\textbf{Generate} \\ ($K$ Responses)};
  \node[datblk, right=of gen, minimum height=2.2cm] (cands) {\textbf{$y^{(1)}$} \\ \textbf{$y^{(2)}$} \\ \textbf{$\vdots$} \\ ($y^{(K)}$)};
  \node[grayblk, right=of cands] (rank) {\textbf{Human/AI} \\ (Rank)};
  \node[lossblk, right=of rank] (loss) {\textbf{Ranking} \\ (Loss)};
  \node[mrgblk, below=0.8cm of loss] (out) {\textbf{Aligned} \\ (Policy)};

  \draw[arr] (x)--(gen);
  \draw[arr] (gen)--(cands);
  \draw[arr] (cands)--(rank);
  \draw[arr] (rank)--(loss);
  \draw[arr] (loss)--(out);
\end{tikzpicture}
\end{adjustbox}
\caption{RRHF: $K$ responses are generated and ranked; policy trained to match rank order via pairwise loss.}
\end{figure}

\subsection{State Transition Table}
\begin{table}[h]\centering
\begin{tabular}{lll}\toprule
\textbf{Property} & \textbf{Before} & \textbf{After}\\\midrule
Responses per prompt & 1 & $K$ (ranked)\\
Supervision signal & Binary (win/lose) & Full rank ordering\\
Loss type & Log-sigmoid & Pairwise margin\\
Top-response training & Separate SFT & Joint SFT + ranking\\
Annotation granularity & Single comparison & Full $K$-way ranking\\\bottomrule
\end{tabular}
\caption{State properties before and after RRHF.}
\end{table}

\newpage
%% ------------------------------------------------------------
\section{Reinforcement Learning from AI Feedback (RLAIF)}

\subsection{Mathematical Foundation}
RLAIF replaces expensive human annotators with a capable AI model (``AI annotator'' $\mathcal{A}$) to generate preference labels at scale. The AI annotator scores pairs via a prompt $p$:
\begin{equation}
  \hat{r}(x,y_i,y_j) = P_\mathcal{A}(y_i \succ y_j \mid p,\, x,\, y_i,\, y_j)
\end{equation}
These AI-generated preferences are then used identically to human preferences to train the RM and run PPO (or DPO):
\begin{equation}
  \mathcal{L}_{\text{RLAIF}} = -\mathbb{E}\!\left[\log\sigma\!\left(r_\phi(x,y_w) - r_\phi(x,y_l)\right)\right], \quad (y_w \succ y_l) \text{ from } \mathcal{A}
\end{equation}

\subsection{Architecture Flow Diagram}

\begin{figure}[h]
\centering
\begin{adjustbox}{max width=\linewidth}
\begin{tikzpicture}[node distance=0.7cm and 1.1cm]
  \node[datblk] (pairs) {\textbf{Response} \\ (Pairs $(y_i, y_j)$)};
  \node[orangeblk, right=of pairs] (ai) {\textbf{AI Annotator} \\ ($\mathcal{A}$ (GPT-4, etc.))};
  \node[datblk, right=of ai] (prefs) {\textbf{AI Preference} \\ (Labels)};
  \node[trnblk, right=of prefs] (rm) {\textbf{Reward} \\ (Model $r_\phi$)};
  \node[mrgblk, below=0.9cm of rm] (ppo) {\textbf{RLHF/DPO} \\ (Training)};
  \node[mrgblk, left=1.4cm of ppo] (out) {\textbf{Aligned} \\ (Policy)};

  \draw[arr] (pairs)--(ai);
  \draw[arr] (ai)--(prefs);
  \draw[arr] (prefs)--(rm);
  \draw[arr] (rm)--(ppo);
  \draw[arr] (ppo)--(out);

  \node[font=\tiny,red,align=center,below=0.3cm of ai] {replaces human\\annotation};
\end{tikzpicture}
\end{adjustbox}
\caption{RLAIF: an AI model generates preference labels at scale, replacing costly human annotators.}
\end{figure}

\subsection{State Transition Table}
\begin{table}[h]\centering
\begin{tabular}{lll}\toprule
\textbf{Property} & \textbf{Before (RLHF)} & \textbf{After (RLAIF)}\\\midrule
Annotation source & Humans & AI annotator $\mathcal{A}$\\
Annotation cost & High (\$\$\$) & Low (API calls)\\
Annotation scale & Limited (thousands) & Massive (millions)\\
Annotation consistency & Variable & Consistent per model\\
Annotation bias & Human bias & AI model bias\\\bottomrule
\end{tabular}
\caption{State properties comparing RLHF and RLAIF.}
\end{table}

\newpage
%% ------------------------------------------------------------
\section{Constitutional AI}

\subsection{Mathematical Foundation}
Constitutional AI uses a set of principles $\mathcal{C} = \{c_1, \ldots, c_M\}$ (the ``constitution'') to iteratively critique and revise model outputs before training. In the RL phase, an AI Feedback model scores responses against the constitution:
\begin{equation}
  r_{\text{CAI}}(x, y) = \frac{1}{M}\sum_{m=1}^{M} P_\mathcal{A}(\text{harmless} \mid c_m, x, y)
\end{equation}
Two training phases: (1) supervised revision via critique-revise loop, (2) RLHF with AI-generated preference labels from the constitution.

\subsection{Architecture Flow Diagram}

\begin{figure}[h]
\centering
\begin{adjustbox}{max width=\linewidth}
\begin{tikzpicture}[node distance=0.65cm and 1.0cm]
  % Critique-revise loop
  \node[datblk] (red) {\textbf{Red-team} \\ (Prompt $x$)};
  \node[trnblk, right=of red] (gen) {\textbf{Generate} \\ (Response $y$)};
  \node[orangeblk, right=of gen] (crit) {\textbf{Critique vs} \\ (Constitution $\mathcal{C}$)};
  \node[trnblk, right=of crit] (rev) {\textbf{Revise} \\ ($y \to y'$)};

  % SFT on revisions
  \node[lossblk, below=1.0cm of rev] (sft) {\textbf{SFT on} \\ (Revisions)};

  % RL phase
  \node[grayblk, below=1.0cm of sft] (aifb) {\textbf{AI Feedback} \\ (Score $r_{\text{CAI}}$)};
  \node[mrgblk, left=2.4cm of aifb] (rl) {\textbf{RLHF/PPO} \\ (Training)};
  \node[mrgblk, left=of rl] (out) {\textbf{Safe} \\ (Policy)};

  \draw[arr] (red)--(gen);
  \draw[arr] (gen)--(crit);
  \draw[arr] (crit)--(rev);
  \draw[arr] (rev.south) -- ++(0,-0.3) -| (gen.south) node[midway,below,font=\tiny]{repeat};
  \draw[arr] (rev)--(sft);
  \draw[arr] (sft)--(aifb);
  \draw[arr] (aifb)--(rl);
  \draw[arr] (rl)--(out);
\end{tikzpicture}
\end{adjustbox}
\caption{Constitutional AI: critique-revise loop creates SFT data; AI feedback from constitution guides RLHF.}
\end{figure}

\subsection{State Transition Table}
\begin{table}[h]\centering
\begin{tabular}{lll}\toprule
\textbf{Property} & \textbf{Before} & \textbf{After}\\\midrule
Safety signal & None / Human labels & AI + constitutionl principles\\
Training phases & SFT only & Critique-revise SFT + RLAIF\\
Red-team coverage & Manual & Systematic via constitution\\
Harmlessness & Moderate & Substantially improved\\
Helpfulness trade-off & N/A & Explicitly balanced\\\bottomrule
\end{tabular}
\caption{State properties before and after Constitutional AI training.}
\end{table}

\newpage
%% ============================================================
\part{Reinforcement Learning Training Methods}
%% ============================================================

\section{Proximal Policy Optimization (PPO)}

\subsection{Mathematical Foundation}
PPO optimizes policy parameters $\theta$ using a clipped objective function to prevent destabilizing parameter updates. Let $\rho_t(\theta) = \frac{\pi_\theta(a_t \mid s_t)}{\pi_{\theta_{\text{old}}}(a_t \mid s_t)}$ be the action probability ratio. The surrogate objective maximizes:
\begin{equation}
  \mathcal{L}^{\text{CLIP}}(\theta) = \hat{\mathbb{E}}_t \left[ \min\left( \rho_t(\theta) \hat{A}_t, \, \text{clip}(\rho_t(\theta), 1-\epsilon, 1+\epsilon) \hat{A}_t \right) \right]
\end{equation}
where $\epsilon$ is a clipping hyperparameter (typically $0.1$ or $0.2$), and $\hat{A}_t$ is the advantage estimated using Generalized Advantage Estimation (GAE):
\begin{equation}
  \hat{A}_t = \sum_{l=0}^\infty (\gamma \lambda)^l \delta_{t+l}^V, \quad \delta_t^V = r_t + \gamma V_\phi(s_{t+1}) - V_\phi(s_t)
\end{equation}
where $V_\phi$ is the value network parameters, $\gamma$ is the discount factor, and $\lambda$ controls bias-variance trade-offs. The value function loss is clipped to match policy updates:
\begin{equation}
  \mathcal{L}^{\text{VF}}(\phi) = \hat{\mathbb{E}}_t \left[ \max\left( (V_\phi(s_t) - V_t^{\text{target}})^2, \, (V_{\phi_{\text{old}}}(s_t) + \text{clip}(V_\phi(s_t) - V_{\phi_{\text{old}}}(s_t), -c, c) - V_t^{\text{target}})^2 \right) \right]
\end{equation}

\subsection{Architecture Flow Diagram}

\begin{figure}[h]
\centering
\begin{adjustbox}{max width=\linewidth}
\begin{tikzpicture}[node distance=0.7cm and 1.1cm]
  \node[datblk, minimum size=1.6cm] (env) {\textbf{Prompt} \\ (Environment)};
  \node[trnblk, right=of env] (actor) {\textbf{Actor} \\ ($\pi_\theta$)};
  \node[trnblk, below=0.8cm of actor] (critic) {\textbf{Critic} \\ ($V_\phi$)};
  \node[grayblk, right=of actor] (traj) {\textbf{Token} \\ (Trajectory)};
  \node[orangeblk, right=of traj] (rm) {\textbf{Reward} \\ ($r_\phi$)};
  \node[lossblk, below right=0.5cm and 1.0cm of traj] (adv) {\textbf{Advantage} \\ ($\hat{A}_t$)};
  \node[mrgblk, right=1.6cm of adv] (upd) {\textbf{PPO} \\ (Update)};

  \draw[arr] (env)--(actor);
  \draw[arr] (actor)--(traj);
  \draw[arr] (traj)--(rm);
  \draw[arr] (traj.south) -| (critic.east);
  \draw[arr] (rm.south) -| (adv.east);
  \draw[arr] (critic.east) -| (adv.west);
  \draw[arr] (adv)--(upd);
  \draw[arr,dashed,blue!60] (upd.east) -- ++(0.4,0) -- ++(0,2.2) -| (actor.north) node[near start,above,font=\tiny]{clip $\rho \in [1\pm\epsilon]$};
\end{tikzpicture}
\end{adjustbox}
\caption{PPO: actor generates trajectory, critic estimates value, advantage drives clipped surrogate update.}
\end{figure}

\subsection{State Transition Table}
\begin{table}[h]\centering
\begin{tabular}{lll}\toprule
\textbf{Property} & \textbf{Before} & \textbf{After}\\\midrule
Policy update size & Unconstrained & Clipped by $\epsilon$\\
Value network & None & Trained $V_\phi$\\
Advantage estimation & N/A & GAE $\hat{A}_t$\\
Training stability & Unstable (vanilla PG) & Stable (clip constraint)\\
Reward signal & Sparse (end of seq) & Discounted returns\\\bottomrule
\end{tabular}
\caption{State properties before and after PPO setup.}
\end{table}

\newpage
%% ------------------------------------------------------------
\section{Group Relative Policy Optimization (GRPO)}

\subsection{Mathematical Foundation}
GRPO optimizes policies by calculating rewards relative to a group of sampled outputs, eliminating the need for a separate critic model.
For a given prompt $x$, the policy generates a group of $G$ candidate completions $\{y_1, y_2, \dots, y_G\}$. Each candidate $y_i$ is evaluated by a reward function (or judge) yielding score $r_i$.
The relative advantage $A_i$ of each candidate is computed as:
\begin{equation}
  A_i = \frac{r_i - \text{mean}(r)}{\text{std}(r)} = \frac{r_i - \frac{1}{G}\sum_j r_j}{\sqrt{\frac{1}{G}\sum_j (r_j - \bar{r})^2 + \epsilon}}
\end{equation}
The policy objective maximizes the clipped advantage with a KL penalty directly evaluated on the generated tokens:
\begin{equation}
  \mathcal{L}_{\text{GRPO}}(\theta) = -\frac{1}{G} \sum_{i=1}^G \left[ \min\left( \rho_i(\theta) A_i, \, \text{clip}(\rho_i(\theta), 1-\epsilon, 1+\epsilon) A_i \right) - \beta D_{\text{KL}}\left( \pi_\theta(y_i \mid x) \parallel \pi_{\text{ref}}(y_i \mid x) \right) \right]
\end{equation}
where the probability ratio $\rho_i(\theta)$ is defined as:
\begin{equation}
  \rho_i(\theta) = \frac{\pi_\theta(y_i \mid x)}{\pi_{\theta_{\text{old}}}(y_i \mid x)}
\end{equation}

\subsection{Architecture Flow Diagram}

\begin{figure}[h]
\centering
\begin{adjustbox}{max width=\linewidth}
\begin{tikzpicture}[node distance=0.6cm and 1.0cm]
  \node[datblk] (x) {\textbf{Prompt $x$}};
  \node[trnblk, right=of x] (pol) {\textbf{Policy $\pi_\theta$}};
  \node[grayblk, right=of pol, minimum height=2.0cm] (group) {\textbf{Group of} \\ \textbf{$G$ Responses} \\ ($y^{(1)}\ldots y^{(G)}$)};
  \node[orangeblk, right=of group] (scores) {\textbf{Rewards} \\ ($r^{(1)}\ldots r^{(G)}$)};
  \node[lossblk, right=of scores] (adv) {\textbf{Normalised} \\ (Advantage $\hat{A}^{(i)}$)};
  \node[mrgblk, below=0.9cm of adv] (upd) {\textbf{GRPO} \\ (Update)};

  \draw[arr] (x)--(pol);
  \draw[arr] (pol)--(group);
  \draw[arr] (group)--(scores);
  \draw[arr] (scores)--(adv);
  \draw[arr] (adv)--(upd);

  \node[red,font=\tiny,align=center,below=0.3cm of pol] {no critic\\network};
\end{tikzpicture}
\end{adjustbox}
\caption{GRPO: group of responses scored together; within-group normalised reward replaces value network.}
\end{figure}

\subsection{State Transition Table}
\begin{table}[h]\centering
\begin{tabular}{lll}\toprule
\textbf{Property} & \textbf{Before (PPO)} & \textbf{After (GRPO)}\\\midrule
Critic / value network & Required & Not needed\\
Advantage estimation & GAE (per token) & Group normalisation\\
Memory overhead & $2\times$ model (actor + critic) & $\approx 1\times$ model\\
Samples per prompt & 1 & $G$ (group sampling)\\
Variance reduction & Value baseline & Mean group reward\\\bottomrule
\end{tabular}
\caption{State properties comparing PPO and GRPO.}
\end{table}

\newpage
%% ------------------------------------------------------------
\section{Rejection Fine-Tuning (RFT)}

\subsection{Mathematical Foundation}
RFT generates $K$ candidate solutions per problem using the current policy, filters by a correctness verifier $\mathcal{V}$, and fine-tunes on the correct subset — essentially a self-improvement loop without RL:
\begin{equation}
  \mathcal{D}_{\text{correct}} = \{(x, y^{(k)}) : \mathcal{V}(x, y^{(k)}) = 1,\; y^{(k)} \sim \pi_\theta,\; k=1..K\}
\end{equation}
\begin{equation}
  \mathcal{L}_{\text{RFT}} = -\sum_{(x,y) \in \mathcal{D}_{\text{correct}}} \sum_t \log P_\theta(y_t \mid x, y_{<t})
\end{equation}
The process repeats: train on correct outputs, generate new candidates, filter again.

\subsection{Architecture Flow Diagram}

\begin{figure}[h]
\centering
\begin{adjustbox}{max width=\linewidth}
\begin{tikzpicture}[node distance=0.7cm and 1.0cm]
  \node[datblk] (prob) {\textbf{Problems} \\ ($\{x_i\}$)};
  \node[trnblk, right=of prob] (pol) {\textbf{Policy $\pi_\theta$} \\ (Sample $K$)};
  \node[grayblk, right=of pol, minimum height=1.8cm] (cands) {\textbf{$K$ Candidate} \\ (Solutions)};
  \node[orangeblk, right=of cands] (ver) {\textbf{Verifier} \\ ($\mathcal{V}$)};
  \node[datblk, right=of ver] (correct) {\textbf{Correct} \\ (Subset)};
  \node[lossblk, below=0.9cm of correct] (sft) {\textbf{SFT on} \\ (Correct)};
  \node[trnblk, left=2.4cm of sft] (newpol) {\textbf{Updated} \\ (Policy)};

  \draw[arr] (prob)--(pol);
  \draw[arr] (pol)--(cands);
  \draw[arr] (cands)--(ver);
  \draw[arr] (ver)--(correct);
  \draw[arr] (correct)--(sft);
  \draw[arr] (sft)--(newpol);
  \draw[arr,dashed,green!60!black] (newpol.north) -| (pol.south) node[midway,below,font=\tiny]{iterate};
\end{tikzpicture}
\end{adjustbox}
\caption{RFT: policy samples $K$ solutions, verifier filters correct ones, SFT improves the policy iteratively.}
\end{figure}

\subsection{State Transition Table}
\begin{table}[h]\centering
\begin{tabular}{lll}\toprule
\textbf{Property} & \textbf{Before} & \textbf{After}\\\midrule
Data source & Static human labels & Self-generated correct solutions\\
RL component & Required (PPO) & None (SFT only)\\
Verifier & None & Required (exact-match, unit tests)\\
Pass@K improvement & Baseline & Improves with iterations\\
Training complexity & High (RL) & Standard SFT loop\\\bottomrule
\end{tabular}
\caption{State properties before and after RFT.}
\end{table}

\newpage
%% ------------------------------------------------------------
\section{REINFORCE}

\subsection{Mathematical Foundation}
REINFORCE is the foundational policy gradient algorithm. For a trajectory $\tau = (s_0,a_0,s_1,a_1,\ldots)$, the gradient of expected return is:
\begin{equation}
  \nabla_\theta J(\theta) = \mathbb{E}_{\tau \sim \pi_\theta}\!\left[G(\tau) \nabla_\theta \log \pi_\theta(\tau)\right],\quad G(\tau) = \sum_{t=0}^T \gamma^t r_t
\end{equation}
For LLMs, the trajectory is the token sequence, each token is an action, and $r_T$ (end reward) is discounted back. A baseline $b$ (e.g., mean reward) reduces variance:
\begin{equation}
  \nabla_\theta J \approx \frac{1}{N}\sum_n (G_n - b)\nabla_\theta \log \pi_\theta(y_n|x_n)
\end{equation}

\subsection{Architecture Flow Diagram}

\begin{figure}[h]
\centering
\begin{adjustbox}{max width=\linewidth}
\begin{tikzpicture}[node distance=0.7cm and 1.0cm]
  \node[datblk] (state) {\textbf{State $s$} \\ (Prompt)};
  \node[trnblk, right=of state] (pol) {\textbf{Policy} \\ ($\pi_\theta$)};
  \node[grayblk, right=of pol] (act) {\textbf{Action $a$} \\ (token)};
  \node[datblk, right=of act] (rew) {\textbf{Reward} \\ ($r_T$)};
  \node[lossblk, below=0.9cm of rew] (ret) {\textbf{Return $G$} \\ ($-$ baseline $b$)};
  \node[mrgblk, left=2.4cm of ret] (upd) {\textbf{PG} \\ (Update $\nabla\log\pi$)};

  \draw[arr] (state)--(pol);
  \draw[arr] (pol)--(act);
  \draw[arr] (act)--(rew);
  \draw[arr] (rew)--(ret);
  \draw[arr] (ret)--(upd);
  \draw[arr,dashed,green!60!black] (upd.north) -| (pol.south) node[midway,below,font=\tiny]{update $\theta$};
\end{tikzpicture}
\end{adjustbox}
\caption{REINFORCE: policy samples action sequence, receives return $G$, updates via policy gradient.}
\end{figure}

\subsection{State Transition Table}
\begin{table}[h]\centering
\begin{tabular}{lll}\toprule
\textbf{Property} & \textbf{Before} & \textbf{After}\\\midrule
Gradient type & Supervised CE & Policy gradient $G\nabla\log\pi$\\
Variance & N/A & High (needs baseline)\\
Critic network & Not needed & Not needed\\
Reward type & N/A & Any differentiable or sparse\\
Update frequency & Per batch & Per episode/trajectory\\\bottomrule
\end{tabular}
\caption{State properties before and after REINFORCE setup.}
\end{table}

\newpage
%% ------------------------------------------------------------
\section{Advantage Actor-Critic (A2C / A3C)}

\subsection{Mathematical Foundation}
A2C uses a shared backbone with two heads: a policy head (actor) and a value head (critic). The advantage $A(s,a) = Q(s,a) - V(s)$ reduces gradient variance:
\begin{equation}
  \mathcal{L}_{\text{actor}} = -\mathbb{E}_t\bigl[A(s_t,a_t)\log\pi_\theta(a_t|s_t)\bigr]
\end{equation}
\begin{equation}
  \mathcal{L}_{\text{critic}} = \mathbb{E}_t\bigl[(r_t + \gamma V(s_{t+1}) - V(s_t))^2\bigr]
\end{equation}
A3C extends A2C with asynchronous parallel workers each running independent environments and gradients are asynchronously applied to a central parameter server.

\subsection{Architecture Flow Diagram}

\begin{figure}[h]
\centering
\begin{adjustbox}{max width=\linewidth}
\begin{tikzpicture}[node distance=0.7cm and 1.0cm]
  \node[frzblk, minimum height=2.0cm] (backbone) {\textbf{Shared} \\ \textbf{Transformer} \\ (Backbone)};
  \node[trnblk, above right=0.5cm and 1.3cm of backbone] (actor) {\textbf{Policy} \\ (Head $\pi_\theta$)};
  \node[trnblk, below right=0.5cm and 1.3cm of backbone] (critic) {\textbf{Value} \\ (Head $V_\phi$)};
  \node[datblk, left=1.5cm of backbone] (env) {\textbf{Environment} \\ ($(s_t, r_t)$)};
  \node[lossblk, right=3.0cm of backbone] (adv) {\textbf{Advantage} \\ ($A(s,a)$)};
  \node[mrgblk, right=of adv] (upd) {\textbf{Joint} \\ (Update)};

  \draw[arr] (env)--(backbone);
  \draw[arr] (backbone.east|-actor)--(actor.west);
  \draw[arr] (backbone.east|-critic)--(critic.west);
  \draw[arr] (actor.east) -| (adv.north);
  \draw[arr] (critic.east) -| (adv.south);
  \draw[arr] (adv)--(upd);
  \draw[arr,dashed,green!60!black] (upd.south) -- ++(0,-0.4) -| (backbone.south);
\end{tikzpicture}
\end{adjustbox}
\caption{A2C: shared backbone feeds actor (policy) and critic (value) heads; advantage combines their outputs.}
\end{figure}

\subsection{State Transition Table}
\begin{table}[h]\centering
\begin{tabular}{lll}\toprule
\textbf{Property} & \textbf{Before (REINFORCE)} & \textbf{After (A2C)}\\\midrule
Variance reduction & Baseline only & Full advantage estimation\\
Value network & Separate & Shared backbone, separate head\\
Gradient stability & Low & Higher (advantage)\\
Parallelism & None & A3C: asynchronous workers\\
Sample efficiency & Low & Higher (on-policy)\\\bottomrule
\end{tabular}
\caption{State properties comparing REINFORCE and A2C.}
\end{table}

\newpage
%% ------------------------------------------------------------
\section{Self-Play Reinforcement Learning}

\subsection{Mathematical Foundation}
In Self-Play RL, the model competes against previous versions of itself. At iteration $n$, the opponent is a snapshot $\pi_{\theta_{n-1}}$ (or mixture of past policies). The reward is determined by competitive outcome:
\begin{equation}
  r(y_{\theta}, y_{\theta_{\text{old}}}) = \begin{cases} +1 & \text{if } y_{\theta} \text{ is judged better} \\ -1 & \text{otherwise} \end{cases}
\end{equation}
\begin{equation}
  \mathcal{J} = \mathbb{E}_{\pi_\theta, \pi_{\theta_{\text{old}}}}\bigl[r(y_\theta, y_{\theta_{\text{old}}})\bigr]
\end{equation}
This creates an ever-escalating difficulty curriculum since the opponent improves alongside the model.

\subsection{Architecture Flow Diagram}

\begin{figure}[h]
\centering
\begin{adjustbox}{max width=\linewidth}
\begin{tikzpicture}[node distance=0.7cm and 1.0cm]
  \node[trnblk] (cur) {\textbf{Current Policy} \\ ($\pi_\theta$ (trains))};
  \node[frzblk, below=1.0cm of cur] (old) {\textbf{Past Policy} \\ ($\pi_{\theta_{\text{old}}}$ (frozen snapshot))};
  \node[datblk, right=2.0cm of $(cur)!0.5!(old)$] (game) {\textbf{Head-to-Head} \\ (Comparison)};
  \node[orangeblk, right=of game] (judge) {\textbf{Judge /} \\ (Verifier $r$)};
  \node[lossblk, right=of judge] (rl) {\textbf{RL Update} \\ ($\nabla\log\pi \cdot r$)};
  \node[grayblk, below=0.9cm of rl] (snap) {\textbf{Snapshot} \\ (Archive)};

  \draw[arr] (cur.east) -| (game.north);
  \draw[arr] (old.east) -| (game.south);
  \draw[arr] (game)--(judge);
  \draw[arr] (judge)--(rl);
  \draw[arr] (rl.south)--(snap.north);
  \draw[arr,dashed,green!60!black] (snap.south) -- ++(0,-0.6) -| (old.south) node[near start,below,font=\tiny]{update snapshot};
  \draw[arr,dashed,green!60!black] (rl.north) -- ++(0,1.5) -| (cur.north) node[near start,above,font=\tiny]{train};
\end{tikzpicture}
\end{adjustbox}
\caption{Self-Play RL: current policy competes against a frozen past snapshot; win/loss drives RL updates.}
\end{figure}

\subsection{State Transition Table}
\begin{table}[h]\centering
\begin{tabular}{lll}\toprule
\textbf{Property} & \textbf{Before} & \textbf{After}\\\midrule
Opponent & None / static & Evolving past self\\
Difficulty curriculum & Fixed & Auto-scaling\\
Reward source & External judge & Comparative win/loss\\
Policy snapshot archive & N/A & Maintained\\
Capability ceiling & Fixed by data & Self-determined\\\bottomrule
\end{tabular}
\caption{State properties before and after Self-Play RL setup.}
\end{table}

\newpage
%% ------------------------------------------------------------
\section{Monte Carlo Tree Search RL (MCTS-RL)}

\subsection{Mathematical Foundation}
MCTS-RL builds a reasoning tree by alternating between \emph{selection} (UCT formula), \emph{expansion}, \emph{rollout}, and \emph{backpropagation}. The UCT score for node $n$:
\begin{equation}
  \text{UCT}(n) = \frac{W(n)}{N(n)} + c\sqrt{\frac{\ln N(\text{parent}(n))}{N(n)}}
\end{equation}
where $W(n)$ is total rollout reward and $N(n)$ is visit count. The policy is improved by training on trajectories weighted by MCTS visit counts:
\begin{equation}
  \pi_\theta \leftarrow \pi_\theta + \eta\,\mathbb{E}\bigl[\text{MCTS visits}(a|s) \cdot \nabla_\theta\log\pi_\theta(a|s)\bigr]
\end{equation}

\subsection{Architecture Flow Diagram}

\begin{figure}[h]
\centering
\begin{adjustbox}{max width=\linewidth}
\begin{tikzpicture}[node distance=0.5cm and 0.9cm]
  \node[datblk,minimum width=1.6cm] (root) {\textbf{Root} \\ (State)};
  \node[trnblk, right=0.7cm of root,minimum width=1.6cm] (sel) {\textbf{Select} \\ (UCT)};
  \node[trnblk, right=0.6cm of sel,minimum width=1.6cm] (exp) {\textbf{Expand} \\ (Node)};
  \node[orangeblk, right=0.6cm of exp,minimum width=1.6cm] (roll) {\textbf{Rollout} \\ (to end)};
  \node[lossblk, right=0.6cm of roll,minimum width=1.6cm] (back) {\textbf{Backup} \\ ($W, N$)};
  \node[mrgblk, below=0.9cm of $(sel)!0.5!(back)$] (pol) {\textbf{Policy Update} \\ (from visit counts)};

  \draw[arr] (root)--(sel);
  \draw[arr] (sel)--(exp);
  \draw[arr] (exp)--(roll);
  \draw[arr] (roll)--(back);
  \draw[arr] (back.south) -- ++(0,-0.3) -| (pol.east);
  \draw[arr] (pol.west) -| (sel.south) node[near start,below,font=\tiny]{next iter};
\end{tikzpicture}
\end{adjustbox}
\caption{MCTS-RL: select (UCT), expand, rollout, backpropagate; visit counts guide policy improvement.}
\end{figure}

\subsection{State Transition Table}
\begin{table}[h]\centering
\begin{tabular}{lll}\toprule
\textbf{Property} & \textbf{Before} & \textbf{After}\\\midrule
Search strategy & Greedy / beam & Tree search (UCT)\\
State value estimate & None & $W(n)/N(n)$ from rollouts\\
Planning horizon & Local & Multi-step look-ahead\\
Compute per query & Inference only & Inference + tree expansion\\
Policy quality & Single-pass & Iteratively refined\\\bottomrule
\end{tabular}
\caption{State properties before and after MCTS-RL integration.}
\end{table}

\newpage
%% ------------------------------------------------------------
\section{Search-Augmented RL (Search-RL)}

\subsection{Mathematical Foundation}
Search-RL augments the reasoning process with an explicit search over candidate next steps. At each step $t$, the policy generates $B$ candidates $\{s_t^{(1)},\ldots,s_t^{(B)}\}$, a verifier scores them, and the best is selected:
\begin{equation}
  s_t^* = \arg\max_{b} V(x, s_{1:t-1}, s_t^{(b)})
\end{equation}
Training uses the full selected trajectory as RL experience:
\begin{equation}
  \mathcal{L}_{\text{SearchRL}} = -\sum_t \mathbb{1}[\mathcal{V}(x,y) = 1]\log\pi_\theta(s_t^* \mid x, s_{1:t-1})
\end{equation}

\subsection{Architecture Flow Diagram}

\begin{figure}[h]
\centering
\begin{adjustbox}{max width=\linewidth}
\begin{tikzpicture}[node distance=0.6cm and 1.0cm]
  \node[datblk] (x) {\textbf{Problem $x$}};
  \node[trnblk, right=of x] (pol) {\textbf{Policy $\pi_\theta$} \\ (beam $B$)};
  \node[grayblk, right=of pol, minimum height=1.8cm] (cands) {\textbf{Step} \\ \textbf{Candidates} \\ ($s^{(1..B)}$)};
  \node[orangeblk, right=of cands] (ver) {\textbf{Step} \\ (Verifier $V$)};
  \node[trnblk, right=of ver] (best) {\textbf{Best Step} \\ ($s^*$)};
  \node[mrgblk, below=0.9cm of best] (train) {\textbf{RL Train} \\ (on trajectory)};

  \draw[arr] (x)--(pol);
  \draw[arr] (pol)--(cands);
  \draw[arr] (cands)--(ver);
  \draw[arr] (ver)--(best);
  \draw[arr] (best.south)--(train.north);
  \draw[arr,dashed,green!60!black] (best.east) -- ++(0.4,0) -- ++(0,1.4) -| (pol.north) node[near start,above,font=\tiny]{next step};
\end{tikzpicture}
\end{adjustbox}
\caption{Search-RL: at each step, beam search generates candidates; verifier selects best; policy trained on trajectory.}
\end{figure}

\subsection{State Transition Table}
\begin{table}[h]\centering
\begin{tabular}{lll}\toprule
\textbf{Property} & \textbf{Before} & \textbf{After}\\\midrule
Step selection & Greedy / top-1 & Verifier-guided beam search\\
Search breadth & 1 & $B$ candidates per step\\
Verifier & None & Step-level verifier $V$\\
Training signal & Outcome reward & Step-selected trajectory\\
Error recovery & None & Search can backtrack\\\bottomrule
\end{tabular}
\caption{State properties before and after Search-RL.}
\end{table}

\newpage
%% ------------------------------------------------------------
\section{Tree-of-Thought Training}

\subsection{Mathematical Foundation}
Tree-of-Thought (ToT) Training teaches the model to generate, evaluate, and prune multiple reasoning branches. The model is trained on annotated trees where each branch $b_i$ has a quality score $q_i \in [0,1]$. The training objective rewards correct branches and penalises dead-ends:
\begin{equation}
  \mathcal{L}_{\text{ToT}} = -\sum_i q_i \sum_t \log P_\theta(b_{i,t} \mid x, b_{i,<t}) + (1-q_i)\sum_t \log(1 - P_\theta(b_{i,t} \mid x, b_{i,<t}))
\end{equation}
A separate scoring model $s_\phi(x, b_i)$ is trained to evaluate intermediate thought nodes and guide beam search at inference.

\subsection{Architecture Flow Diagram}

\begin{figure}[h]
\centering
\begin{adjustbox}{max width=\linewidth}
\begin{tikzpicture}[node distance=0.5cm and 0.8cm]
  \node[datblk,minimum width=1.6cm] (root) {\textbf{Root} \\ (Thought)};
  \node[trnblk,minimum width=1.4cm, above right=0.3cm and 0.8cm of root] (b1) {\textbf{Branch 1} \\ (good)};
  \node[lossblk,minimum width=1.4cm, right=0.3cm and 0.8cm of root] (b2) {\textbf{Branch 2} \\ (bad)};
  \node[trnblk,minimum width=1.4cm, below right=0.3cm and 0.8cm of root] (b3) {\textbf{Branch 3} \\ (good)};
  \node[mrgblk,minimum width=1.4cm, right=0.9cm of b1] (ans1) {\textbf{Ans $A_1$}};
  \node[mrgblk,minimum width=1.4cm, right=0.9cm of b3] (ans3) {\textbf{Ans $A_3$}};
  \node[grayblk, below=1.0cm of b3] (score) {\textbf{Branch} \\ (Scorer $s_\phi$)};

  \draw[arr] (root)--(b1);
  \draw[arr] (root)--(b2);
  \draw[arr] (root)--(b3);
  \draw[arr] (b1)--(ans1);
  \draw[arr] (b3)--(ans3);
  \draw[arr,red,dashed] (b2) -- ++(0.6,0) node[right,font=\tiny,red]{pruned};
  \draw[arr] (b1.west) -- ++(-0.2,0) |- (score.west);
  \draw[arr] (b2.west) -- ++(-0.4,0) |- (score.west);
  \draw[arr] (b3.south) -- (score.north);
\end{tikzpicture}
\end{adjustbox}
\caption{ToT Training: good branches lead to answers; bad branches are pruned; scorer trained to evaluate all.}
\end{figure}

\subsection{State Transition Table}
\begin{table}[h]\centering
\begin{tabular}{lll}\toprule
\textbf{Property} & \textbf{Before} & \textbf{After}\\\midrule
Reasoning structure & Linear chain & Branching tree\\
Dead-end handling & None & Explicit pruning\\
Branch evaluator & None & Trained $s_\phi$\\
Search at inference & Linear & Guided tree beam\\
Complex planning & Weak & Strong\\\bottomrule
\end{tabular}
\caption{State properties before and after Tree-of-Thought Training.}
\end{table}

\newpage
%% ------------------------------------------------------------
\section{Self-Consistency Training}

\subsection{Mathematical Foundation}
Self-Consistency Training generates $K$ independent reasoning paths and uses majority voting to determine the correct answer. The model is then trained to increase probability of the majority-vote answer $y^*$:
\begin{equation}
  y^* = \arg\max_{y} \sum_{k=1}^K \mathbf{1}[\text{answer}(y^{(k)}) = y],\quad y^{(k)} \sim \pi_\theta(\cdot|x)
\end{equation}
\begin{equation}
  \mathcal{L}_{\text{SC}} = -\sum_{\{k:\,\text{answer}(y^{(k)})=y^*\}} \sum_t \log P_\theta(y^{(k)}_t \mid x, y^{(k)}_{<t})
\end{equation}
Paths that agree with the majority vote are reinforced; minority paths are not trained on (or penalised).

\subsection{Architecture Flow Diagram}

\begin{figure}[h]
\centering
\begin{adjustbox}{max width=\linewidth}
\begin{tikzpicture}[node distance=0.6cm and 1.0cm]
  \node[datblk] (x) {\textbf{Prompt $x$}};
  \node[trnblk, right=of x] (pol) {\textbf{Sample $K$} \\ (Paths)};
  \node[grayblk, right=of pol, minimum height=2.0cm] (paths) {\textbf{$y^{(1)}\ldots y^{(K)}$} \\ (different chains)};
  \node[orangeblk, right=of paths] (vote) {\textbf{Majority} \\ (Vote $y^*$)};
  \node[trnblk, right=of vote] (winners) {\textbf{Winning} \\ (Paths)};
  \node[lossblk, below=0.9cm of winners] (sft) {\textbf{SFT on} \\ (winners)};
  \node[mrgblk, left=1.6cm of sft] (out) {\textbf{Updated} \\ ($\pi_\theta$)};

  \draw[arr] (x)--(pol);
  \draw[arr] (pol)--(paths);
  \draw[arr] (paths)--(vote);
  \draw[arr] (vote)--(winners);
  \draw[arr] (winners)--(sft);
  \draw[arr] (sft)--(out);
\end{tikzpicture}
\end{adjustbox}
\caption{Self-Consistency Training: majority-vote over $K$ paths identifies winning chains; model trained on those.}
\end{figure}

\subsection{State Transition Table}
\begin{table}[h]\centering
\begin{tabular}{lll}\toprule
\textbf{Property} & \textbf{Before} & \textbf{After}\\\midrule
Inference mode & Single greedy & $K$ diverse samples\\
Answer selection & Argmax & Majority vote\\
Training target & All completions & Majority-consistent paths only\\
Accuracy ceiling & Limited by single pass & Higher (ensemble effect)\\
Compute cost & $1\times$ & $K\times$ at sampling time\\\bottomrule
\end{tabular}
\caption{State properties before and after Self-Consistency Training.}
\end{table}

\newpage
%% ============================================================
\part{Reward Modeling}
%% ============================================================

\section{Reward Model (RM)}

\subsection{Mathematical Foundation}
A Reward Model takes a prompt $x$ and a completion $y$ and outputs a scalar preference score. It is trained on human-ranked pairs $(y_w \succ y_l)$ via the Bradley-Terry log-likelihood:
\begin{equation}
  \mathcal{L}_{\text{RM}} = -\mathbb{E}_{(x,y_w,y_l)}\bigl[\log\sigma(r_\phi(x,y_w) - r_\phi(x,y_l))\bigr]
\end{equation}
The RM is typically initialised from a fine-tuned LLM with the final token's hidden state projected to a scalar via a linear head $W_r \in \mathbb{R}^{d \times 1}$:
\begin{equation}
  r_\phi(x,y) = W_r^\top h_{\text{last}}(x,y)
\end{equation}

\subsection{Architecture Flow Diagram}

\begin{figure}[h]
\centering
\begin{adjustbox}{max width=\linewidth}
\begin{tikzpicture}[node distance=0.7cm and 1.1cm]
  \node[datblk] (inp) {\textbf{$(x, y)$} \\ (Input)};
  \node[frzblk, right=of inp] (backbone) {\textbf{LLM} \\ (Backbone)};
  \node[grayblk, right=of backbone] (last) {\textbf{Last Token} \\ (Hidden $h$)};
  \node[trnblk, right=of last] (head) {\textbf{Scalar} \\ (Head $W_r$)};
  \node[mrgblk, right=of head] (score) {\textbf{Reward} \\ (Score $r$)};

  \node[datblk, above=1.0cm of backbone] (yw) {\textbf{Chosen $y_w$}};
  \node[datblk, below=1.0cm of backbone] (yl) {\textbf{Rejected $y_l$}};
  \node[lossblk, below=1.0cm of head] (loss) {\textbf{BT Loss} \\ ($\log\sigma(r_w - r_l)$)};

  \draw[arr] (inp)--(backbone);
  \draw[arr] (backbone)--(last);
  \draw[arr] (last)--(head);
  \draw[arr] (head)--(score);
  \draw[arr] (yw.south)--(backbone.north);
  \draw[arr] (yl.north)--(backbone.south);
  \draw[arr] (score.south) |- (loss.east);
\end{tikzpicture}
\end{adjustbox}
\caption{Reward Model: LLM backbone maps $(x,y)$ pair to scalar via linear head; trained on chosen/rejected pairs.}
\end{figure}

\subsection{State Transition Table}
\begin{table}[h]\centering
\begin{tabular}{lll}\toprule
\textbf{Property} & \textbf{Before} & \textbf{After}\\\midrule
Output type & Token distribution & Scalar reward\\
Final layer & LM head & Scalar head $W_r$\\
Training data & Text completions & $(x, y_w, y_l)$ preference pairs\\
Loss function & Cross-entropy & Bradley-Terry log-likelihood\\
Usage & Language generation & RLHF reward signal\\\bottomrule
\end{tabular}
\caption{State properties before and after reward model training.}
\end{table}

\newpage
%% ------------------------------------------------------------
\section{Process Reward Model (PRM)}

\subsection{Mathematical Foundation}
A PRM assigns a reward to each intermediate reasoning step $s_t$ rather than to the final answer alone. Given a solution prefix $(s_1,\ldots,s_t)$, the PRM outputs a binary step quality $q_t \in \{0,1\}$:
\begin{equation}
  r_{\text{PRM}}(x, s_{1:T}) = \prod_{t=1}^T r_t,\quad r_t = P_\phi(q_t = 1 \mid x, s_{1:t})
\end{equation}
The PRM is trained on step-level human labels with binary cross-entropy:
\begin{equation}
  \mathcal{L}_{\text{PRM}} = -\sum_t \bigl[q_t \log r_t + (1-q_t)\log(1-r_t)\bigr]
\end{equation}

\subsection{Architecture Flow Diagram}

\begin{figure}[h]
\centering
\begin{adjustbox}{max width=\linewidth}
\begin{tikzpicture}[node distance=0.5cm and 0.8cm]
  \node[datblk,minimum width=1.6cm] (x) {\textbf{Problem} \\ ($x$)};
  \node[trnblk,minimum width=1.4cm, right=0.7cm of x] (s1) {\textbf{Step 1} \\ ($s_1$)};
  \node[trnblk,minimum width=1.4cm, right=0.6cm of s1] (s2) {\textbf{Step 2} \\ ($s_2$)};
  \node[grayblk,minimum width=0.6cm, right=0.6cm of s2] (dots) {\textbf{$\cdots$}};
  \node[trnblk,minimum width=1.4cm, right=0.6cm of dots] (sT) {\textbf{Step $T$} \\ ($s_T$)};

  \node[orangeblk,minimum width=1.4cm, below=0.7cm of s1] (r1) {\textbf{$r_1$}};
  \node[orangeblk,minimum width=1.4cm, below=0.7cm of s2] (r2) {\textbf{$r_2$}};
  \node[orangeblk,minimum width=1.4cm, below=0.7cm of sT] (rT) {\textbf{$r_T$}};
  \node[mrgblk, right=0.6cm of rT] (prod) {\textbf{$\prod r_t$} \\ (PRM score)};

  \foreach \s/\r in {s1/r1, s2/r2, sT/rT}
    \draw[arr] (\s)--(\r);
  \draw[arr] (r1.east) -- (r2.west);
  \draw[arr] (r2.east) -- ++(0.3,0);
  \draw[arr] (rT)--(prod);
  \foreach \a/\b in {x/s1, s1/s2, s2/dots, dots/sT}
    \draw[arr] (\a)--(\b);
\end{tikzpicture}
\end{adjustbox}
\caption{PRM: each reasoning step independently scored; product of step scores gives overall process reward.}
\end{figure}

\subsection{State Transition Table}
\begin{table}[h]\centering
\begin{tabular}{lll}\toprule
\textbf{Property} & \textbf{Before (RM)} & \textbf{After (PRM)}\\\midrule
Reward granularity & Final answer only & Per-step\\
Error localisation & None & Identifies failing step\\
Label annotation & Answer-level & Step-level (expensive)\\
Reward signal density & Sparse (1 per traj.) & Dense ($T$ per traj.)\\
Training on correct solutions & Required & Not required\\\bottomrule
\end{tabular}
\caption{State properties comparing RM and PRM.}
\end{table}

\newpage
%% ------------------------------------------------------------
\section{Outcome Reward Model (ORM)}

\subsection{Mathematical Foundation}
An ORM judges only the \emph{correctness of the final answer} $y_{\text{ans}}$, ignoring intermediate steps. For verifiable tasks (math, code), correctness can be checked automatically:
\begin{equation}
  r_{\text{ORM}}(x, y) = \mathbf{1}[\text{verify}(y_{\text{ans}}, y^*_{\text{gold}})]
\end{equation}
For open-ended tasks, the ORM is a learned binary classifier trained on (correct, incorrect) outcome labels:
\begin{equation}
  \mathcal{L}_{\text{ORM}} = -\bigl[c\log P_\phi(y) + (1-c)\log(1-P_\phi(y))\bigr],\quad c \in \{0,1\}
\end{equation}

\subsection{Architecture Flow Diagram}

\begin{figure}[h]
\centering
\begin{adjustbox}{max width=\linewidth}
\begin{tikzpicture}[node distance=0.7cm and 1.1cm]
  \node[datblk] (x) {\textbf{Problem $x$}};
  \node[trnblk, right=of x] (sol) {\textbf{Full} \\ (Solution $y$)};
  \node[frzblk, right=of sol] (lm) {\textbf{ORM} \\ ($P_\phi$)};
  \node[lossblk, right=of lm] (check) {\textbf{Outcome} \\ (Check)};
  \node[mrgblk, right=of check] (score) {\textbf{Binary} \\ (Score $c$)};

  \node[datblk,below=0.8cm of lm] (gold) {\textbf{Gold Answer} \\ ($y^*$)};

  \draw[arr] (x)--(sol);
  \draw[arr] (sol)--(lm);
  \draw[arr] (lm)--(check);
  \draw[arr] (gold.north)-|(check.south);
  \draw[arr] (check)--(score);

  \node[font=\tiny,green!60!black,align=center,above=0.2cm of check] {exact-match\\or learned};
\end{tikzpicture}
\end{adjustbox}
\caption{ORM: scores the final answer against gold; can be rule-based (exact-match) or learned classifier.}
\end{figure}

\subsection{State Transition Table}
\begin{table}[h]\centering
\begin{tabular}{lll}\toprule
\textbf{Property} & \textbf{Before} & \textbf{After}\\\midrule
Reward scope & N/A & Final answer only\\
Intermediate credit & None & None\\
Annotation cost & N/A & Low (binary label)\\
Verifiable tasks & Manual check & Automatic (rule-based)\\
Open-ended tasks & None & Learned classifier\\\bottomrule
\end{tabular}
\caption{State properties before and after ORM setup.}
\end{table}

\newpage
%% ------------------------------------------------------------
\section{Verifier Model}

\subsection{Mathematical Foundation}
A Verifier Model is a binary classifier that determines whether a candidate solution is correct. Unlike an ORM that outputs a probability, verifiers typically use a hard decision rule learned from (correct, incorrect) solution pairs:
\begin{equation}
  \hat{v}(x, y) = \mathbf{1}\bigl[P_\phi(\text{correct} \mid x, y) \geq \tau\bigr]
\end{equation}
Verifiers are used in test-time search (Best-of-N) and training-time filtering (RFT). The decision threshold $\tau$ trades off precision and recall on the verification task:
\begin{equation}
  \mathcal{L}_{\text{ver}} = -\bigl[v\log P_\phi + (1-v)\log(1-P_\phi)\bigr],\quad v \in \{0,1\}
\end{equation}

\subsection{Architecture Flow Diagram}

\begin{figure}[h]
\centering
\begin{adjustbox}{max width=\linewidth}
\begin{tikzpicture}[node distance=0.7cm and 1.1cm]
  \node[datblk] (x) {\textbf{$(x, y)$} \\ (Solution)};
  \node[frzblk, right=of x] (enc) {\textbf{LLM} \\ (Encoder)};
  \node[trnblk, right=of enc] (cls) {\textbf{Binary} \\ (Classifier $P_\phi$)};
  \node[grayblk, right=of cls] (thr) {\textbf{Threshold} \\ ($\tau$)};
  \node[trnblk, above right=0.4cm and 0.9cm of thr] (yes) {\textbf{Correct $\checkmark$}};
  \node[lossblk, below right=0.4cm and 0.9cm of thr] (no) {\textbf{Incorrect $\times$}};

  \draw[arr] (x)--(enc);
  \draw[arr] (enc)--(cls);
  \draw[arr] (cls)--(thr);
  \draw[arr] (thr.east) |- (yes.west);
  \draw[arr] (thr.east) |- (no.west);
\end{tikzpicture}
\end{adjustbox}
\caption{Verifier Model: encodes solution, classifies as correct/incorrect via learned threshold.}
\end{figure}

\subsection{State Transition Table}
\begin{table}[h]\centering
\begin{tabular}{lll}\toprule
\textbf{Property} & \textbf{Before} & \textbf{After}\\\midrule
Output & Continuous score & Binary decision\\
Threshold & N/A & Tunable $\tau$\\
Primary use & Ranking & Filtering (accept/reject)\\
Training data & Ranked pairs & Correct/incorrect labels\\
Precision/recall trade-off & N/A & Controlled via $\tau$\\\bottomrule
\end{tabular}
\caption{State properties before and after Verifier Model training.}
\end{table}

\newpage
%% ------------------------------------------------------------
\section{Critic Model}

\subsection{Mathematical Foundation}
A Critic Model generates natural-language \emph{critiques} of candidate outputs and optionally a numerical score. It is trained on (output, critique, revision) triplets. The critique generation loss:
\begin{equation}
  \mathcal{L}_{\text{crit}} = -\sum_t \log P_\phi(c_t \mid x, y, c_{<t})
\end{equation}
where $c$ is the gold critique. Critics enable self-refinement loops: the model generates $y$, the critic evaluates it producing $c$, and the model revises $y \to y'$ conditioned on $(x, y, c)$:
\begin{equation}
  y' \sim P_\theta(\cdot \mid x, y, c)
\end{equation}

\subsection{Architecture Flow Diagram}

\begin{figure}[h]
\centering
\begin{adjustbox}{max width=\linewidth}
\begin{tikzpicture}[node distance=0.7cm and 1.0cm]
  \node[datblk] (x) {\textbf{Prompt $x$}};
  \node[trnblk, right=of x] (gen) {\textbf{Generator} \\ ($\pi_\theta$)};
  \node[grayblk, right=of gen] (y) {\textbf{Draft $y$}};
  \node[orangeblk, right=of y] (crit) {\textbf{Critic} \\ ($P_\phi$)};
  \node[grayblk, right=of crit] (c) {\textbf{Critique $c$} \\ (+ Score)};
  \node[trnblk, below=0.9cm of c] (rev) {\textbf{Revise} \\ ($y \to y'$)};
  \node[mrgblk, left=2.0cm of rev] (out) {\textbf{Final} \\ (Output $y'$)};

  \draw[arr] (x)--(gen);
  \draw[arr] (gen)--(y);
  \draw[arr] (y)--(crit);
  \draw[arr] (crit)--(c);
  \draw[arr] (c)--(rev);
  \draw[arr] (y.south) |- (rev.east);
  \draw[arr] (rev)--(out);
  \draw[arr,dashed,green!60!black] (rev.north) -- ++(0,0.5) -| (gen.north) node[midway,above,font=\tiny]{optional loop};
\end{tikzpicture}
\end{adjustbox}
\caption{Critic Model: generates natural-language critique of a draft; generator revises using the critique.}
\end{figure}

\subsection{State Transition Table}
\begin{table}[h]\centering
\begin{tabular}{lll}\toprule
\textbf{Property} & \textbf{Before} & \textbf{After}\\\midrule
Feedback type & Scalar reward & Natural-language critique\\
Interpretability & Low (scalar) & High (textual feedback)\\
Revision loop & Not applicable & Generator refines on critique\\
Training data & Preference pairs & (output, critique, revision) triplets\\
Use case & RM scoring & Self-refinement / RLHF\\\bottomrule
\end{tabular}
\caption{State properties before and after Critic Model training.}
\end{table}

